\section{A falsifiable program: where exactness should and should not help}
\label{sec:program}

The two probes at controlled scale (Sections~\ref{sec:maze}
and~\ref{sec:deq}) resolved to attainability, not inductive bias: the
pre-registered matched-fit design found the exact solve wins because the
alternatives cannot fit the task, not because it tips training toward a
better solution among equally-attainable ones. What survives is the
operator's classical, verified numerical properties -- exact, deterministic,
$\bigO(\text{fill})$ -- against incumbents that are stochastic, iterative, or
conditioning-dependent. This section states that operator gap, the
benchmark-scale experiments that would show it matters (or does not), the
negative control, and the single prerequisite that gates all of them. None of
what follows claims inductive bias; each hypothesis below is a numerical or
capability claim, stated so it can be wrong.

\paragraph{The operator gap (what scalable differentiable-sparse methods must match).}
Three capabilities anchor the claim; each is exact and deterministic where the
incumbents are stochastic, iterative, or conditioning-dependent.
\begin{itemize}
  \item \textbf{M1: exact, deterministic} gradient of $\log\det A$ and exact
        on-pattern blocks and diagonal of $\inv{A}$, at $\bigO(\text{fill})$.
        Incumbent: the conjugate-gradient plus stochastic-Lanczos-quadrature
        gradient of GPyTorch \cite{gardner2018gpytorch,ubaru2017fast} is
        variance-laden and $\kappa$-biased, with no exact $\diagop(\inv{A})$
        at scale.
  \item \textbf{M2: conditioning-independent, constant-memory adjoint.} A
        direct factorization gives an $\bigO(\text{fill})$-memory,
        iteration-free adjoint whose accuracy does not degrade with $\kappa$
        or $\rho(J)$ (Section~\ref{sec:deq}). Incumbents: unrolled iterative
        adjoints store $\bigO(K)$ states; implicit CG and AMG adjoints
        reintroduce $\kappa$-dependent iteration and gradient bias.
  \item \textbf{M3: a one-layer, exact, differentiable global graph
        functional} (resolvent, effective resistance, commute time), where
        bounded-depth message passing provably needs $\Omega(\text{diameter})$
        depth and transformers approximate without extrapolating
        (Section~\ref{sec:maze}).
\end{itemize}

\paragraph{Benchmark-scale experiments.}
Each attacks one capability on real, non-toy data, with a falsification bar
that should be fixed before the benchmark is run.
\begin{itemize}
  \item \textbf{E2 (M1): exact deterministic gradients for calibrated UQ.}
        Hypothesis: exact marginal-likelihood gradients and exact posterior
        marginal variances give better-calibrated uncertainty and more stable
        hyperparameter optimization than CG plus SLQ, with the gap widening as
        the noise tends to zero and $\kappa$ grows. Competitors: GPyTorch
        \cite{gardner2018gpytorch}, SKI/KISS-GP \cite{wilson2015kernel},
        Vecchia approximations \cite{katzfuss2021general}, and SVGP
        \cite{hensman2013gaussian}. Data: ERA5/NOAA spatial GMRFs and PeMS
        state-space GPs. \emph{Falsified if} SLQ/CG gradients are within noise
        of exact on real data at feasible cost.
  \item \textbf{E1 (M2): constant-memory, conditioning-independent adjoints
        for differentiable physics.} Hypothesis: on high-contrast 2-D
        PDE-constrained learning, the exact sparse-direct adjoint gives
        gradients whose error and backward memory are flat in the condition
        number, where unrolled CG runs out of memory and implicit AMG adjoints
        stall or bias. Competitors: unrolled CG and implicit CG/AMG adjoints.
        Data: high-contrast Darcy flow on SPE10 permeability (contrast up to
        $10^7$); a sparse-observation inverse problem. \emph{Falsified if}
        AMG-preconditioned implicit adjoints stay accurate and cheap at $10^7$
        contrast.
  \item \textbf{E3 (M3): the exact resolvent layer vs.\ over-squashing and
        size extrapolation.} Hypothesis: one exact differentiable resolvent
        layer computes global functionals that bounded-depth message passing
        cannot, removing over-squashing
        \cite{alon2021bottleneck,topping2022understanding} and generalizing
        across graph sizes where GNNs, rewiring, and graph transformers
        degrade. This is the benchmark-scale version of
        Section~\ref{sec:maze}. Competitors: strong MPNNs, graph rewiring
        (SDRF \cite{topping2022understanding}, FoSR \cite{karhadkar2023fosr}),
        graph transformers (GPS \cite{rampasek2022recipe}, Exphormer
        \cite{shirzad2023exphormer}), and implicit GNNs \cite{gu2020implicit}.
        Data: the Long Range Graph Benchmark \cite{dwivedi2022long} and real
        road networks with train-small/test-large splits. \emph{Falsified if}
        a graph transformer matches both in distribution and across the
        extrapolation sweep.
  \item \textbf{E4 (M2, deepest): training at the edge of stability.} The
        benchmark-scale version of Section~\ref{sec:deq}: does the exact
        backward let one train structured equilibrium models at near-critical
        $\rho(J)\to1$, a regime that phantom-gradient \cite{geng2021training}
        and Anderson/Broyden adjoints
        \cite{anderson1965iterative,broyden1965class} cannot reach, for a
        learned-capability gain (longer effective range)? \emph{Falsified if} phantom-gradient
        DEQs train fine without ever needing high $\rho$, i.e.\ the useful
        regime is not stiff.
\end{itemize}

\paragraph{The negative control.}
Every generalization experiment includes a task \emph{not} reducible to a
structured global solve, where any exactness-specific explanation,
attainability included, predicts the advantage
must vanish. A persistent gap there would mean the effect is a confound
(capacity, optimization, data leakage), not exactness; its absence is what
supports the attribution reading of Section~\ref{sec:maze} at scale.

Two pre-registered negative controls have already been run at small scale,
and neither passed cleanly, which is exactly why Sections~\ref{sec:maze}
and~\ref{sec:deq} resolve to the conservative attainability reading rather
than a confirmed inductive bias. On the maze held-out seeds, a
short-correlation-length variant of the routing task ($\varepsilon=2.0$,
where the solve's reach should barely matter) still showed the exact solve
beating the compute-matched truncation at every size -- a disclosed limitation
of the matched-compute cost story. On the DEQ held-out seeds, the low-$\rho$
silent control's relative-gap clause failed at near-machine-precision loss
levels, though its independent function-distance clause held by more than
four orders of magnitude (Section~\ref{sec:deq}). Per the frozen rule, a
failed negative control voids a clean positive reading; neither probe gets
one. The benchmark-scale negative controls above carry the same commitment:
if they fail, the failure is reported, and a future benchmark-scale
inductive-bias test would need a cleaner control (keyed on function distance,
not relative gap, at near-machine-precision cells) before it could confirm
anything.

\paragraph{The one prerequisite, and the one falsification that sinks it.}
All of the above need problems larger than today's per-node Python-loop CPU
ceiling. The gating item is a single, already-scoped engineering build, a
GPU-batched junction kernel, not new theory; until it runs, the claim is
established only at the controlled scale of Sections~\ref{sec:maze}
and~\ref{sec:deq}. One outcome would bound the whole program: if real
low-treewidth problems are mostly well-conditioned (where iterative and
stochastic gradients are good enough) and real stiff problems are mostly dense
3-D (multigrid territory), then the operator gap lands where it is not needed.
E1 and E2 are designed so that their failure pins down that
boundary, in which case the result is a characterization of where exactness
matters.
